\section{Scope and Positioning}

This paper sits at the intersection of uncertainty-aware recognition, grouped decoding, and low-resource glyph and handwriting recognition. It is intentionally positioned as a bounded empirical study of uncertainty propagation rather than as a general recognition framework.

Two scope boundaries matter throughout. First, the paper separates fully real evidence from synthetic-from-real evidence. The symbol-level evidence in the preserved workshop package is real. The Omniglot, Digits, and Kuzushiji-49 sequence and downstream tasks are synthetic-from-real in structure. The grouped Historical Newspapers and ScaDS.AI experiments are real grouped evaluations, and the redesigned downstream task is derived from those real grouped transcripts. Second, the paper is not a decipherment or semantic-recovery study. It does not claim plaintext recovery, manuscript understanding, or a universal law of uncertainty propagation.

This positioning also clarifies what is novel here. The strongest contribution is not another decoder comparison in isolation. Instead, the paper contributes a support-aware account of when rescued alternatives survive long enough to matter at grouped and downstream levels, and when they do not.
